\section{$N=2$ 超对称规范理论的 Seiberg--Witten 分析}

\begin{keypoints}
\begin{itemize}[leftmargin=1.4em,itemsep=2pt,topsep=2pt,parsep=0pt]
  \item \textbf{物理图像：}$N=2$ 纯 $SU(2)$ 理论在模空间（坐标 $u=\langle{\rm Tr}\,\phi^2\rangle$）上处于库仑相；远离奇点时低能自由度只有未破缺的 $U(1)^{r}$ 矢量多重态，BPS 界 $M\geq\sqrt{2}\,|Z|$ 与中心荷 $Z=a\,n_e+a_D\,n_m$ 是全节的出发点。
  \item \textbf{数学工具：}低能有效作用量由全纯预势 ${\cal F}(A)$ 完全决定（式 (3.4)--(3.5)），$\tau(a)=\partial^2{\cal F}/\partial a^2$ 同时给出模空间度规；单圈预势 (3.7) 在微扰论中精确，非微扰修正按 $(\Lambda/A)^{4k}$ 展开。
  \item \textbf{关键结论：}手征反常把 $U(1)_{\cal R}$ 破缺为 ${\bf Z}_{4N_c}$（式 (3.6) 与零模计数一致），希格斯真空又把它破缺到 ${\bf Z}_4$（$SU(2)$ 情形）；${\rm Im}\,\tau$ 调和且必须为正，故 $\tau(a)$ 只能局部定义，须用对偶描述拼接。
  \item \textbf{关键结果：}电磁对偶把 $\tau\to-1/\tau$，与 $\tau\to\tau+1$ 一起生成 $SL(2,{\rm Z})$，在 $(a_D,a)$ 上分式线性地作用；为保持双子质量公式 (3.12)，电荷矢量 $(n_m,n_e)$ 必须按逆矩阵 $M^{-1}$ 变换，且 $M$ 为整系数矩阵。
\end{itemize}
\end{keypoints}

本节中，我们按照 Seiberg 与 Witten 的做法 \cite{SW-I}，分析 $N=2$ 纯 $SU(2)$ 理论。$N=2$ 理论具有强有力的 Ward 恒等式，它们与一些物理输入相结合，可以引出诸如磁单极凝聚这样的有意义结论，并且在把 $N=2$ 软破缺到 $N=1$ 之后，引出由无质量磁单极凝聚所导致的禁闭。这条推理链漫长而精细。对于所考虑的这些理论，当把 $\theta$ 角取为零时，BPS 界由下式给出：
\begin{eqnarray}
M &\geq& \sqrt{2}\,\vert Z \vert\,,\nonumber \\
Z &=& a\,(n_e + {i\over \alpha} n_m)\,,\quad{\rm with}\quad 
\alpha = {g^2\over 4\pi}\,. \nonumber
\end{eqnarray}
$Z$ 的形式在同时作 $n_e \leftrightarrow n_m$、$\alpha\leftrightarrow 1/\alpha$ 与 $a \leftrightarrow a/\alpha$ 的交换下保持不变。这可以看作电磁对偶的 Montonen--Olive 猜想的一个佐证，尽管我们已经看到，在 $N=2$ 理论中 $Z$ 会获得一种有意义的重整化。该猜想以其原始形式只在 $N=4$ 理论中才有意义；而对于 $N=2$，我们仍然可以有一个 $SL(2, Z)$ 变换作用在参数 $\tau = \theta/2\pi + i4\pi/g^2$ 上，它体现了 Montonen--Olive 猜想的一种较弱形式。Seiberg 与 Witten 通过把这些变换与整体对称性以及动能正定性的要求结合起来，得以构造出一套确定低能理论精确形式的程序，这将在下面介绍。

% [Tutor Note: 上式中 $Z=a(n_e+\frac{i}{\alpha}n_m)$ 在 $n_e\leftrightarrow n_m$、$\alpha\leftrightarrow1/\alpha$、$a\leftrightarrow a/\alpha$ 下只差一个整体相位（$Z\to iZ$），故 $|Z|$ 不变；这正是 $g\leftrightarrow g_D=4\pi/g$ 的 Montonen--Olive 对偶的“弱形式”。注意这只是经典（树图）层面的观察：$N=2$ 量子理论中 $Z$ 的重整化通过 $a_D=\partial{\cal F}/\partial a$ 进入，精确公式是 (3.12)。]

% [Tutor Note: $\tau=\theta/2\pi+i4\pi/g^2$ 的定义极易记错：虚部是 $4\pi/g^2$（不是 $1/g^2$，也不是 $g^2/4\pi$），因此 $\operatorname{Im}\tau=4\pi/g^2>0$ 恰对应弱耦合极限 $g^2\to0$，即 $\tau$ 位于上半平面（upper half plane）。与本节导语的 $\alpha=g^2/4\pi$ 对照有 $\tau=\theta/2\pi+i/\alpha$。还要区分：此处的 $\tau$ 是经典（树图）耦合；在量子模空间上 $\tau(a)=\partial^2{\cal F}/\partial a^2$（见式 (3.5)），两者仅在 $|a|\to\infty$ 的渐近自由区重合。]

\subsection{模空间的参数化}

$N=2$ 超对称 Yang--Mills 理论 (\ref{N=2}) 的经典作用量中含有标量势

% [Tutor Note: 原文 “containts” 系 “contains” 的拼写笔误（原文 PDF 第 43 页）。势能前的因子 $1/2g^2$ 与 $\operatorname{Tr}$ 的归一化约定绑定（见 §2 的约定）；若改用基本表示下的 $\operatorname{tr}$，前置因子会随之改变。$V\geq0$，且 $V=0\iff[\phi,\phi^\dagger]=0$，这正是下文定义希格斯真空的依据。]
$$
V={1\over 2g^2}{\rm Tr}\left([\phi^\dagger,\phi]^2\right)\,.
$$
由此可见，希格斯真空由 $[\phi,\phi^\dagger]=0$ 定义，这意味着 $\phi$ 取值于规范群的嘉当（Cartan）子代数：$\phi= \phi_i H^i$。于是，一般地，规范群 $G$ 破缺到由嘉当子代数中元素生成的子群 $H$。$G/H$ 中的元素并不使希格斯真空保持不变，但作为规范变换，它们联系着物理上等价的真空。另一方面，一旦选定了嘉当子代数的某一组基，则该子代数内 $\phi$ 的不同真空值就对应于不同的物理理论。因此，$\phi$ 中的这些自由度（即各 $\phi_i$）参数化了物理上不等价真空构成的空间，也就是理论的模空间。该模空间的维数等于规范群 $G$ 的秩 $r$。然而，这一模空间参数化并非我们所需要的形式，因为仍然存在残余的规范不变性：陪集 $G/H$ 中含有这样的元素，它们虽不使真空保持不变，却不会把 $\phi$ 带出嘉当子代数。这些变换正是 Weyl 反射。因此，模空间的正确参数化不是由 $\phi$ 给出，而是由 $\phi$ 构造出的 Weyl 不变函数给出。

% [Tutor Note: 这里的逻辑链值得补全：$G/H$ 中保持嘉当子代数不变的元素构成 Weyl 群，它通过共轭作用把 $\phi$ 在嘉当内的坐标作反射；物理真空由 Weyl 轨道分类，故模空间的坐标函数必须是 Weyl 不变量。Chevalley 定理进一步保证：Weyl 不变多项式环是自由多项式环，恰有 $r$ 个生成元，这与“模空间维数 = 秩 $r$”一致。原文只说“Weyl 反射”而未点明这层不变量理论，容易让读者误以为 $G/H$ 全是规范冗余。]

Weyl 不变量由特征方程
$$
\det (\lambda - \phi ) = 0\,.
$$
由于 Weyl 反射通过共轭作用作用于 $\phi$，$\det (\lambda-\phi)$ 是不变的。因此，若把该量按 $\lambda$ 的幂展开成多项式，则其系数都是 Weyl 不变量。下面我们把这些系数表示为 $\phi$ 的简单函数。设 $a_1,\cdots , a_N$ 表示特征方程的根，亦即 $\phi$ 的本征值。若规范群为 $SU(N)$（或 $SO(N)$），则 ${\rm Tr}\phi=\sum a_i=0$。对本征值 $a_i$ 的一般取值，规范群破缺到 $U(1)^r$。当某些本征值重合时，未破缺的群从 $U(1)^r$ 跳跃为至少与 $U(1)^{r-1}\times SU(2)$ 一样大的群。用 $a_i$ 表示时，特征多项式具有如下形式：
$$
\lambda^N + \lambda^{N-2} \sum_{i<j} a_ia_j - \lambda^{N-3}
\sum_{i<j<k} a_ia_ja_k + \cdots + (-1)^N \prod_{i=1}^N a_i = 0\,.
$$
对 $SU(2)$，$\phi={1\over 2}a\sigma_3$，容易验证所需的 Weyl 不变量为 $u={\rm Tr}(\phi^2)={1\over2}a^2$。对 $SU(3)$，特征多项式中的系数是 $a_1a_2+a_1a_3+a_2a_3$ 与 $a_1 a_2 a_3$，其中 $a_1, a_2, a_3$ 是 $\phi$ 的本征值。利用 $({\rm Tr}\phi)^2=0$ 与 $({\rm Tr}\phi)^3=0$，容易把 Weyl 不变量写成

% [Tutor Note: “容易验证”一步的细节（原文 PDF 第 44 页）：$\phi={1\over2}a\sigma_3$ 时 $\phi^2={1\over4}a^2\mathbb{1}$，在基本表示中 $\operatorname{Tr}\mathbb{1}=2$，故 ${\rm Tr}(\phi^2)={1\over2}a^2$；Weyl 反射 $a\to-a$ 使之不变，而秩为 1 说明这是唯一的独立不变量。]

% [Tutor Note: $({\rm Tr}\phi)^2=0$、$({\rm Tr}\phi)^3=0$ 是缩写，应理解为把 ${\rm Tr}\phi=\sum_i a_i$ 代入后恒为零所得的牛顿恒等式：由 $(\sum_i a_i)^2=0$ 得 $\sum_{i<j}a_ia_j=-{1\over2}\sum_i a_i^2$，于是 $u={1\over2}\sum_i a_i^2=-\sum_{i<j}a_ia_j$；由 $(\sum_i a_i)^3=0$ 并利用三变量对称性可得 $a_1a_2a_3={1\over3}\sum_i a_i^3$。按字面读会把这两个恒等式误解为空话。]
$$
\begin{array}{ll}
u=\frac{1}{2}{\rm Tr}(\phi^2)=-(a_1a_2+a_1a_3+a_2a_3)\,,  \\
v=-\frac{1}{3}{\rm Tr}(\phi^3)=a_1 a_2 a_3\,.
\end{array}
$$
% [Tutor Note: 原文此处有一个符号笔误（原文 PDF 第 44 页）：与上一页的一般展开 $\det(\lambda-\phi)=\lambda^N-\frac12{\rm Tr}(\phi^2)\lambda^{N-2}-\frac13{\rm Tr}(\phi^3)\lambda^{N-3}+\cdots$ 对照，$N=3$ 时常数项为 $-\frac13{\rm Tr}(\phi^3)=-(a_1a_2a_3)$，故正确的定义应是 $v=+\frac13{\rm Tr}(\phi^3)=a_1a_2a_3$；原文写成 $v=-\frac13{\rm Tr}(\phi^3)=a_1a_2a_3$，中间表达式的负号有误（右端 $a_1a_2a_3$ 才是 $v$ 的定义）。公式已按原文照录，此处如实指出。]

一般地，对 $SU(N)$ 可以算出类似的公式。特征多项式的系数是 $\phi$ 的“Chern”类，
$$
\det(\lambda-\phi)=\lambda^N -\lambda^{N-1} c_{1}(\phi)+
\lambda^{N-2}c_{2}(\phi)+\cdots +(-1)^{j} \lambda^{N-j} c_{j} (\phi )+
\cdots + (-1)^N c_N(\phi)= 0\,.
$$
系数 $c_j(\phi)$ 可由下述形式展开容易地确定：
\begin{eqnarray}
\det(\lambda -\phi) &=& \lambda^N \det (1-\phi/\lambda) =
\lambda^N e^{{\rm Tr\,ln}(1-\phi/\lambda)} 
=\lambda^N {\rm exp}\left(-\sum_{n=1}^\infty \frac{{\rm Tr}\,(\phi^n)}
{n\lambda^n}\right) \nonumber \\
&=& \lambda^N - \frac{1}{2}{\rm Tr}(\phi^2)\lambda^{N-2} -\frac{1}{3}
{\rm Tr}(\phi^3) \lambda^{N-3}+\cdots\,. \nonumber
\end{eqnarray} 
注意，${\rm ln} (1-\phi/\lambda)$ 的级数展开只在 $\lambda >> \phi$ 时才有意义。因此，在上述展开中只有 $\lambda$ 的正幂次项是相关的，它们提供了我们为模空间的规范不变参数化所需的全部 Weyl 不变量。

上面我们把 $\phi$ 当作经典量处理。在量子场论中，我们用相应经典 Weyl 不变量的真空期望值来参数化模空间。例如，$SU(2)\,\,N=2$ 超对称 Yang--Mills 理论的模空间由 $u=<{\rm Tr}(\phi^2)>$ 参数化，在经典极限下它化为 $a^2/2$。

% [Tutor Note: 原文用 $<$ 与 $>$ 表示真空期望值（$u=<{\rm Tr}(\phi^2)>$），这是老式 LaTeX 写法，排版上与不等号同形；规范记法应为 $\langle{\rm Tr}(\phi^2)\rangle$。数学内容未改动，仅在此提示。]

\subsection{$R$ 对称性及其破缺}

$N$ 扩展超对称代数 (\ref{susy-noC}) 在 $N$ 个超荷的整体 $U(N)$ 旋转下不变。因此，一个超对称理论也应当表现出这种整体对称性，通常称之为 $R$ 对称性。这一整体对称性对超荷的作用，可以容易地翻译成超空间变量的变换。例如，对 $N=1$，$U(1)$ $R$ 对称性以 $Q\rightarrow e^{-i\alpha}Q$ 作用于超荷。由此可以得到它对超空间坐标的作用为 $\theta\rightarrow e^{i\alpha}\theta$ 与 $\bar\theta\rightarrow e^{-i\alpha}\bar\theta$（见 (\ref{superspace-Q})）。对 $N=2$，我们有一个如上作用于 $N=2$ 超空间坐标 $\theta^I, \bar\theta_I$ 的 $U(1)$，以及一个旋转超荷指标 $I$ 的 $SU(2)$ $R$ 对称性。为使超对称拉格朗日量在这些变换下不变，我们必须给各种超场赋予适当的变换性质。由此可以容易地得到各分量场在 $R$ 对称性下的行为。下面我们将稍为详细地说明这一点。

% [Tutor Note: 这一步“容易翻译”的实质：超荷的微分表示中含 $\partial/\partial\theta$ 项，要使 $\{Q,\bar Q\}\sim\partial_\mu$ 与 $Q$ 作用于 $\theta$ 的形式不变，$\theta$ 必须与 $Q$ 反向变换；同时 $d^4\theta$ 测度在 $\theta\to e^{i\alpha}\theta$ 下不变，而 $d^2\theta$ 测度带荷 $-2$。这既解释了 (3.1) 中 $\Phi$ 前的因子 $e^{2i\alpha}$（$\Phi$ 的 $U(1)_{\cal R}$ 荷为 2），也解释了超势 ${\cal W}=\sqrt{2}\widetilde Q\Phi Q$ 的总荷必须为 2（见 (3.1) 之后关于 ${\cal W}$ 的讨论）。]

\noindent\underline{$N=2$ 矢量多重态上的作用：}
如上一节所讨论的，$N=2$ 矢量多重态包含一个矢量场 $A_\mu$、两个 Weyl 旋量 $\lambda, \psi$ 和一个标量 $\phi$，它们都在规范群 $G$ 的伴随表示中变换。这些分量可以排成
$$
\begin{array}{ccc}
 & A_\mu & \\
 \lambda & & \psi \\
 & \phi &
\end{array}\,.
$$
$SU(2)_R$ 变换作用在上述图示的各行上，把费米子 $(\lambda,\psi)$ 相互旋转，同时保持 $A_\mu$ 与 $\phi$ 不变。在 $N=1$ 形式中，该多重态分解为一个矢量超场 $V(A_\mu, \lambda)$ 与一个手征超场 $\Phi(\phi,\psi)$。因此，$SU(2)_R$ 中在 $N=1$ 语言下仍保持显式的部分，只是一个不混合 $\lambda$ 与 $\psi$ 的 $U(1)_J$ 子群。该子群由 $\sigma_3$ 生成，作用为 $(\lambda,\psi)\rightarrow(e^{i\alpha}\lambda,e^{-i\alpha}\psi)$。$U(1)_{\cal R}$ 对 $N=1$ 超场的作用已在上一节讨论过。下面我们把这些变换总结如下：

% [Tutor Note: 注意区分 $U(1)_J$ 与 $U(1)_{\cal R}$：在 (3.1) 中 $V$ 的变换完全相同，区别仅在于手征超场 $\Phi$ 是否带整体相位 $e^{2i\alpha}$。$U(1)_J\subset SU(2)_R$ 由 $\sigma_3$ 生成，是 $N=1$ 语言下仍显式的那个 $R$ 对称性；$U(1)_{\cal R}$ 则是 $N=2$ 的额外整体 $U(1)$。两者是独立参数，切勿把它们混成一个 $\alpha$（原文在 (3.2) 中同样用 $\alpha$ 记两个不同 $U(1)$ 的参数）。]
\begin{equation}
\begin{array}{lll}
U(1)_{\cal R} &: \Phi\rightarrow e^{2i\alpha}\Phi(e^{-i\alpha}\theta),
&V\rightarrow V(e^{-i\alpha}\theta)\,,
\\ \\
U(1)_J &: \Phi \rightarrow \Phi (e^{-i\alpha}\theta), &
V\rightarrow V(e^{-i\alpha}\theta)\,.
\end{array}
\label{Rb1}
\end{equation}

\noindent\underline{$N=2$ 超多重态上的作用：}
在 $N=2$ 中，物质场以超多重态的形式出现，每个超多重态包含两个复标量 $(q,\widetilde q^\dagger)$ 与两个 Weyl 费米子 $(\psi_q,\widetilde\psi^\dagger_q)$。所有这些分量都在规范群的同一个（通常是基本）表示中变换。超多重态的各分量可以排成
$$
\begin{array}{ccc}
 & \psi_q & \\
 q & & \widetilde q^\dagger \\
 & \widetilde\psi_q^\dagger &
\end{array}\,. 
$$
同样，$SU(2)_R$ 作用于各行，从而旋转 $q,\widetilde q^\dagger$。在 $N=1$ 语言中，超多重态按手征多重态 $Q(q,\psi_q)$ 与 $\widetilde Q(\widetilde q,\widetilde\psi_q)$ 分解，它们带有互为对偶的规范量子数。在这一分解中，$SU(2)_R$ 只有 $U(1)_J$ 子群是显式的。此外，在 $N=1$ 分解中，与矢量多重态相互作用的超多重态会产生一项超势
$$
{\cal W} = \sqrt{2} \widetilde Q \Phi Q.
$$
由于 ${\cal W}$ 应带有两个单位的 $U(1)_{\cal R}$ 荷，$Q$ 与 $\widetilde Q$ 都是中性的。下面我们总结这些变换：

% [Tutor Note: 由 (3.1) 中 $\Phi\to e^{2i\alpha}\Phi$ 可知 ${\cal W}=\sqrt{2}\widetilde Q\Phi Q$ 的 $U(1)_{\cal R}$ 荷为 $\Phi$ 的荷 2；要使超势不变，$Q,\widetilde Q$ 必须中性——这正是 (3.2) 中它们在 $U(1)_{\cal R}$ 下只随超空间坐标平移而无整体相位的理由。]
\begin{equation}
\begin{array}{lll}
U(1)_{\cal R}: &Q \rightarrow Q(e^{-i\alpha}\theta)\,,
& \widetilde Q \rightarrow \widetilde Q (e^{-i\alpha}\theta)\,, \\ \\
U(1)_J : &Q\rightarrow e^{i\alpha} Q(e^{-i\alpha}\theta)\,,
&\widetilde Q\rightarrow e^{i\alpha}\widetilde Q(e^{-i\alpha}\theta)\,. 
\end{array}
\label{Rb2}
\end{equation}
为后文方便，下面列出所有分量场在 $U(1)_{\cal R}$ 与 $U(1)_J$ 下的变换：
$$
\begin{array}{lccc} 
U(1)_{\cal R}:\qquad & \phi &\rightarrow & e^{2i\alpha}\phi\,, \\
&(\psi,\lambda)&\rightarrow & e^{i\alpha}(\psi,\lambda)\,, \\
&(\psi_q,\widetilde\psi_q)&\rightarrow &e^{-i\alpha} 
(\psi_q,\widetilde\psi_q)\,, \\
&(A_\mu, q, \widetilde q)&\rightarrow & (A_\mu, q, \widetilde q)\,.
\\ \\
U(1)_J: &(\lambda, q)&\rightarrow & e^{i\alpha}(\lambda, q)\,,\\
&(\psi,\widetilde q^\dagger)&\rightarrow & e^{-i\alpha}
(\psi,\widetilde q^\dagger)\,,\\
&(\phi, A_\mu, \psi_q,\widetilde\psi_q^\dagger)&\rightarrow &
(\phi, A_\mu, \psi_q,\widetilde\psi_q^\dagger)\,.
\end{array}
$$
注意，我们可以把二分量旋量 $\lambda$ 与 $\bar\psi$ 组合成一个四分量 Dirac 旋量 $\psi_D$（见 (\ref{sp-comp})）。旋量 $\psi_D$ 在 $U(1)_J$ 下按 $e^{i\alpha}\psi_D$ 变换，在 $U(1)_{\cal R}$ 下按 $e^{i\alpha\gamma_5}\psi_D$ 变换。类似地，$U(1)_{\cal R}$ 对由 $\psi_q$ 与 $\widetilde\psi_q^\dagger$ 构造的 Dirac 旋量表现为一个手征 $U(1)$，只是此时的荷符号相反。因此，$U(1)_{\cal R}$ 是一个手征对称性，从而会被手征反常所破缺，这一点将在下面讨论。

% [Tutor Note: $\psi_D$ 的 $U(1)_{\cal R}$ 变换 $e^{i\alpha\gamma_5}$ 是轴（手征）转动，会在路径积分测度中产生反常，这是 $U(1)_{\cal R}$ 破缺的根源；而 $U(1)_J$ 给出普通的矢量相位 $e^{i\alpha}$（$\lambda$ 与 $\psi$ 反向变换，使 $\psi_D=(\lambda,\bar\psi)$ 获得同一相位），因而是非反常的。两种 $U(1)$ 对手征反常的贡献截然不同，这是本节最易混淆之处。]

\noindent\underline{$R$ 对称性的破缺：}
经典地，我们的理论以完整的整体对称群 $SU(2)_R\times U(1)_{\cal R}$ 为对称群。然而在量子层面，$U(1)_{\cal R}$ 由于反常而破缺到一个离散子群。该子群可以容易地由初等的瞬子考虑确定（更直接的方法将在下一小节给出）。为计算规范群 $SU(N_c)$ 的反常，注意由指标定理，在存在一个瞬子时，基本或反基本表示中的每个左行费米子有一个零模，而伴随表示中的每个左手费米子有 $2N_c$ 个零模。该理论中的关联函数涉及对这些零模所对应的费米子集体坐标的积分。要使关联函数非零，它必须含有足够多的费米子插入以“吃掉”这些零模。因此，在存在 $N_f$ 个味时，第一个非零关联函数是

% [Tutor Note: 指标定理给出的零模数核对：对 $SU(N_c)$，表示 $R$ 中一个左手 Weyl 费米子在 $\nu$ 个瞬子背景中有 $2T(R)\nu$ 个零模；$T({\rm adj})=N_c$（故 $2N_c$ 个），$T({\rm fund})={1\over2}$（故 1 个），与原文一致。原文此处 “left moving”（左行）与 “left-handed”（左手）混用，实指同一件事（手征零模）；“representaion” 是 “representation” 的拼写笔误（原文 PDF 第 46 页）。用集体坐标语言说：瞬子振幅的费米子积分要恰好插入全部零模才非零，下式 (3.3) 就是能“吃掉”零模的最小关联函数，其 $U(1)_{\cal R}$ 荷决定破缺后的离散群。]
\begin{equation}
G=\langle\lambda(x_1)\cdots\lambda(x_{2N_c})\psi(y_1)\cdots\psi(y_{2N_c})
\psi_q(z_1)\cdots\psi_q(z_{N_f})\widetilde\psi_q(u_1)\cdots
\widetilde\psi_q(u_{N_f})\rangle \,.
\label{G}
\end{equation}
在 $U(1)_{\cal R}$ 下，$G$ 按如下方式变换：
$$
G \rightarrow e^{i\alpha(4N_c-2N_f)}\,G\,.
$$
因此 $U(1)_{\cal R}$ 破缺到离散群 ${\bf Z}_{4N_c-2N_f}$。下面我们集中讨论纯 Yang--Mills 理论，即 $N_f=0$。此时 $U(1)_{\cal R}\rightarrow {\bf Z}_{4N_c}$，并以 $e^{2\pi i\alpha}$ 表示，其中 $\alpha=n/4N_c\,, n=1,\cdots 4N_c$。于是整体对称群为 $SU(2)_R\times {\bf Z}_{4N_c}$。然而要注意，$SU(2)_R$ 的中心，其作用为 $(\lambda,\psi)\rightarrow e^{i\pi}(\lambda,\psi)$，也包含在 ${\bf Z}_{4N_c}$ 之中（对应于 $n=2N_c$）。因此，真正的对称群是

% [Tutor Note: 荷的核对（极易错）：由 (3.1) 下方的分量表，$\lambda,\psi$ 的 $U(1)_{\cal R}$ 荷各为 $+1$，$\psi_q,\widetilde\psi_q$ 各为 $-1$，故 (3.3) 的总荷为 $2N_c\cdot1+2N_c\cdot1-N_f\cdot1-N_f\cdot1=4N_c-2N_f$。令上式的相位为 1 得 $\alpha=2\pi n/(4N_c-2N_f)$。注意原文在离散情形把群元素写成 $e^{2\pi i\alpha}$、$\alpha=n/4N_c$，这与前文连续转动参数（荷 $q$ 的场取 $e^{iq\alpha}$）的归一化不同：离散情形下荷 $q$ 的场实际获得 $(e^{2\pi in/(4N_c)})^q$。下文说 $\phi^2$（荷 4）变换为 $e^{2\pi in/N_c}\phi^2$ 即按此约定，可自行核对。]

% [Tutor Note: $n=2N_c$ 对应的群元素为 $e^{2\pi i\cdot2N_c/(4N_c)}=e^{i\pi}=-1$，作用在荷 1 的费米子上给 $-1$，恰与 $SU(2)_R$ 中心 $\lambda,\psi\to e^{i\pi}(\lambda,\psi)$ 相同，故必须取商 ${\bf Z}_2$。这是下文 $SU(2)$ 最终群 $(SU(2)_R\times{\bf Z}_4)/{\bf Z}_2$ 的来源。]
$$
\left(SU(2)_R\times {\bf Z}_{4N_c}\right)/{\bf Z}_2\,.
$$

这一残余的 $R$ 对称性还会被希格斯真空进一步破缺。场 $\phi^2$ 在 ${\bf Z}_{4N_c}$ 下的荷为 4，变换为 $e^{2\pi i n/N_c}\phi^2$，它仅在 $n=N_c,2N_c,3N_c,4N_c$ 时不变。因此，若真空由非零的 $\phi^2$ 标志，则 ${\bf Z}_{4N_c}$ 破缺为 ${\bf Z}_4$。这正是 $SU(2)$ 规范理论（$N_c=2$）的情形。此时，那些不保持 $\phi^2$ 不变的元素都表现为一个 ${\bf Z_2}: \phi^2\rightarrow -\phi^2$。因此，$SU(2)$ 规范理论最终的 $R$ 对称群是 $(SU(2)_R\times {\bf Z_4})/{\bf Z_2}$。

对规范群 $SU(3)$，我们有 ${\bf Z_{4N_c}}={\bf Z_{12}}$。真空由 $\phi^2$ 与 $\phi^3$ 参数化，它们在 ${\bf Z}_{4N_c}$ 下分别变换为 $e^{2\pi i n/N_c}\phi^2$ 与 $e^{3\pi i n/N_c}\phi^3$。$\phi^2$ 的不变性要求 $n=N_c,2N_c,3N_c,4N_c$，使 ${\bf Z_{12}}$ 破缺为 ${\bf Z_{4}}$。$\phi^3$ 的不变性挑出 $n=2N_c, 4N_c$。因此，对 $SU(3)$ 规范理论，${\bf Z}_{12}$ 破缺为 ${\bf Z}_2$。破缺掉的 ${\bf Z}_6$ 子群在模空间上的作用非平凡。对 $SU(4)$ 及更高的规范群，${\bf Z}_{4N_c}$ 不再有任何子群留存下来。

% [Tutor Note: “$\phi^3$ 的不变性挑出 $n=2N_c,4N_c$”是压缩表述（原文 PDF 第 47 页），容易误读：单看 $\phi^3$ 的不变性（相位 $e^{3\pi i n/N_c}=1$）只给出 $n$ 为偶数，即 6 个元素（${\bf Z}_6$，与下一句“破缺掉的 ${\bf Z}_6$”呼应）；必须与 $\phi^2$ 的条件 $n\in\{N_c,2N_c,3N_c,4N_c\}=\{3,6,9,12\}$ 取交集，才剩下 $n=6,12$，即 ${\bf Z}_2$。]

% [Tutor Note: 对 $SU(4)$ 及以上“没有任何子群存活”一句若按字面理解会与计算矛盾（原文 PDF 第 47 页）：元素 $n=2N_c$ 对所有 Weyl 不变量 ${\rm Tr}\phi^k$（$U(1)_{\cal R}$ 荷 $q=2k$，均为偶数）给出的相位为 $e^{2\pi i q\cdot2N_c/(4N_c)}=e^{i\pi q}=1$，即它保持整个真空模空间不变。该元素恰好是 $SU(2)_R$ 的中心（见上一段），在 $(SU(2)_R\times{\bf Z}_{4N_c})/{\bf Z}_2$ 的商下与群中心重合，因而不给出新的对称性。原文的意思应理解为“商掉中心后不再有新的离散 $R$ 对称性”；这一点原文没有点明，特此提醒。]

\subsection{$N=2$ 规范理论的低能有效作用量}

让我们考虑基于秩为 $r$ 的群 $G$ 的 $N=2$ 超对称 Yang--Mills 理论，它被非零的 $\langle\phi\rangle$ 自发破缺。在 $\langle\phi\rangle$ 的两个本征值重合的点之外，唯一的无质量场是与 $G$ 的未破缺子群 $U(1)^{r}$ 相联系的矢量超多重态。在足够低的能量下，没有任何有质量态会作为物理态出现，理论的有效描述可以仅用无质量场给出。原则上，这一低能理论可以通过把低能截断以上的所有有质量模式以及无质量模式积分掉而得到。实践中，这一程序并不容易实施。Seiberg 与 Witten 提出了一个间接的程序来确定精确的低能理论，我们将在本节剩余部分予以说明。

% [Tutor Note: 原文 “non of the massive states” 应为 “none of the massive states”；“as we will in the remaining part of this section” 语句不完整（疑脱漏 “do/describe”）。两处均为原文笔误（原文 PDF 第 47 页），译文按语义补足，未增删物理内容。]
$$
\frac{1}{4\pi} \, {\rm Im} \, \left [ \int \, d^{4}\theta \,
\frac{\partial {\cal F}}{\partial A^{i}} \bar{A}^{i} +
\int \, d^{2} \theta \, \frac{1}{2}
\frac{\partial^{2}{\cal F}}{\partial A_{i} \partial A_{j}}
W^{\alpha i} W^j_\alpha
\right ]\,.
$$
用 ${\cal F}(A)$ 表示时，凯勒势由 $K={\rm Im}\left(\bar{A}_{i}\partial{\cal F}(A)/\partial A_{i}\right)$ 给出。若以 $a_i$ 表示手征超场 $A_i$ 的标量分量，则场空间上的度规（因而也是希格斯真空空间上的度规）由下式给出：

% [Tutor Note: 指标位置极易写反：$K={\rm Im}(\bar A_i\,\partial{\cal F}/\partial A_i)$，度规 $g_{i\bar j}={\rm Im}\,\partial^2{\cal F}/\partial a_i\partial a_j$；注意 $\bar A_i$ 是复共轭并带下指标。另注意归一化：由 (3.4) 的 $1/4\pi$，(3.5) 直接给出 $\tau=\partial^2{\cal F}/\partial a^2$（不再带 $1/4\pi$）；若漏掉 (3.4) 中的 $1/4\pi$，$\tau$ 会与 $\tau_{cl}=\theta/2\pi+i4\pi/g^2$ 相差 $4\pi$，后文所有对偶变换都会随之错位。]
$$
ds^2 = g_{i\bar{j}}\,da^i d\bar{a}^j =
{\rm Im}\,\frac{d^2{\cal F}}{\partial a_i \partial a_j}
\, da^i d\bar{a}^j \,.
$$
注意，凭借 $N=2$ 超对称性，该度规与上述拉格朗日量中 $F_{\mu\nu}^iF^{\mu\nu i}$ 项前出现的广义规范耦合相同。这一关系在量子理论中不会被修正。

% [Tutor Note: “度规 = 广义规范耦合”是 $N=2$ 的核心结构：同一个全纯函数 ${\cal F}$ 同时给出标量场空间的凯勒度规与规范动力学因子（$\operatorname{Im}\tau\,F_{\mu\nu}F^{\mu\nu}$ 的系数）。由于 $N=2$ 的非重整化性质，这一等同关系不因量子修正而改变；这正是后文能用 ${\rm Im}\,\tau>0$（动能正定）约束 $\tau$、并把求度规归结为求 $a_D(u),a(u)$ 的原因。]

下面我们主要关心自发破缺到 $U(1)$ 的 $SU(2)$ 规范群。此时低能有效拉格朗日量具有如下形式
\begin{equation}
{\cal L}_{eff}=
\frac{1}{4\pi} \, {\rm Im} \, \left [ \int \, d^{4}\theta \,
\frac{\partial {\cal F}(A)}{\partial A} \bar{A} +
\int \, d^{2} \theta \, \frac{1}{2}
\frac{\partial^{2}{\cal F}}{\partial A^{2}} W^{\alpha} W_{\alpha}
\right ]\,, 
\label{EA}
\end{equation}
而场空间上的度规由下式给出：
\begin{equation}
ds^2= {\rm Im}(\tau)\, da d\bar a\,,\quad {\rm where},\quad
\tau (a) = \frac{\partial^{2} {\cal F}}{\partial a^{2}}.
\label{metric}
\end{equation}
本小节的目标是仿照 \cite{Seiberg88} 确定 ${\cal F}$ 的微扰形式。确定包含非微扰修正在内的预势的精确形式，是后面几个小节的主题。

为确定 ${\cal F}$ 的单圈微扰修正，我们从规范群为 $SU(2)$ 的精确微观理论 (\ref{N=2}) 出发。该理论是渐近自由的，因此在高能下可以作可靠的微扰计算。如上一小节所指出的，该理论的 $U(1)_{\cal R}$ 对称性被标准的手征反常所破缺。于是我们有

% [Tutor Note: 记号陷阱：此处的 $\alpha$ 是 $U(1)_{\cal R}$ 的转动参数；本节导语中的 $\alpha=g^2/4\pi$；(3.10) 前后 $SL(2,{\rm Z})$ 矩阵的四个元素也记作 $\alpha,\beta,\gamma,\delta$。同一个字母在本节中有三种完全不同的含义，阅读时必须按上下文区分，这是原文记号最易误导之处。]
$$
\partial_\mu J^\mu_5= -\frac{N_c}{8\pi^2}F_{\mu\nu}\widetilde
F^{\mu\nu} \,.
$$
这意味着，在单圈层次上，在一个 $U(1)_{\cal R}$ 变换下，有效拉格朗日量的改变为

% [Tutor Note: 反常系数的归一化易错：原文用 $(32\pi^2)^{-1}\int F\widetilde F\in{\bf Z}$（分母是 $32\pi^2$ 而非 $16\pi^2$），这与 $\theta$ 项的写法 $\theta\,(32\pi^2)^{-1}\int F\widetilde F$ 配套（可由 (3.4)--(3.5) 中 $\theta$ 项的系数验证）。核对离散子群时必须使用同一约定，并记住这里的 $\alpha$ 是 $U(1)_{\cal R}$ 的转角。]
\begin{equation}
\delta {\cal L}_{eff}= - \frac{\alpha N_c}{8\pi^2} F\widetilde F\,.
\label{anom}
\end{equation}
若理论中还含有 $N_f$ 个处于基本表示的费米子，则上述表达式中须把 $N_c$ 替换为 $N_c-N_f/2$。注意，由于 $(32\pi^2)^{-1}\int F\widetilde F$ 是整数，该反常把 $U(1)_{\cal R}$ 破缺为 ${\bf Z}_{4N_c}$（在存在物质时为 ${\bf Z}_{4N_c-2N_f}$）。上一小节从瞬子背景中费米子零模的考虑得到了同样的结果。此外，该反常也可视为引起了 $\theta$ 角的平移。下面我们利用这一自由度，通过对费米子作适当的手征转动把 $\theta$ 取为零。

% [Tutor Note: $N_c\to N_c-N_f/2$ 的来源是零模计数：伴随表示（$\lambda,\psi$，荷各 $+1$）贡献 $2N_c+2N_c=4N_c$，每个基本表示味（$\psi_q,\widetilde\psi_q$，荷各 $-1$）贡献 2，故有效荷为 $4N_c-2N_f=4(N_c-N_f/2)$，离散群为 ${\bf Z}_{4(N_c-N_f/2)}={\bf Z}_{4N_c-2N_f}$。]

% [Tutor Note: 手征转动把 $\theta$ 移走的机制：$U(1)_{\cal R}$ 转动使 $\delta{\cal L}_{eff}=-\alpha N_c/(8\pi^2)F\widetilde F$，而 $\theta$ 项为 $\theta\,(32\pi^2)^{-1}F\widetilde F$，故等价于 $\theta\to\theta-4\pi N_c\alpha$（纯 Yang--Mills，$N_f=0$）；选取适当的 $\alpha$ 即可把 $\theta$ 归零。此后本节默认 $\theta=0$，从而 $\tau_{cl}=i4\pi/g^2$ 为纯虚数。]

${\cal F}$ 的单圈形式可由如下要求确定：在 $U(1)_{\cal R}$ 变换下，低能有效作用量 (\ref{EA}) 的改变应如 (\ref{anom}) 所示。变分 $\delta{\cal L}_{eff}$ 只可能来自 ${\cal L}_{eff}$ 中对 $F_{\mu\nu}$ 为二次的项。只写出拉格朗日量 (\ref{EA}) 中的相关项，这意味着
$$ 
\frac{1}{16\pi}{\rm Im}\left[{\cal F}''(e^{2i\alpha}A)(-FF + iF
\widetilde F)\right] = \frac{1}{16\pi}{\rm Im}\left[{\cal F}''(A)
(-FF+iF\widetilde F)\right]-\frac{\alpha N_c}{8\pi^2}F\widetilde F\,.
$$ 
因此，预势的形式受到如下关系的限制：
$$
{\cal F}''(e^{2i\alpha}A)={\cal F}''(A) - \frac{2\alpha N_c}{\pi}\,,
$$
或者，对无穷小的 $\alpha$，受到如下关系的限制：
$$
\frac{\partial^3 {\cal F}}{\partial A^3}=\frac{N_c}{\pi}\frac{i}{A}\,.
$$
上式可以容易地积分，给出预势的单圈表达式

% [Tutor Note: “容易积分”的细节（原文 PDF 第 49 页）：由 $\partial^3{\cal F}/\partial A^3=(N_c/\pi)(i/A)$ 作两次不定积分得 $\partial^2{\cal F}/\partial A^2=(iN_c/\pi)\ln A+c_1$；$N_c=2$ 时与 $\operatorname{Im}\tau$ 的单圈形式 $\tau={\cal F}''=(i/\pi)(\ln(A^2/\Lambda^2)+3)$ 比较可见，$-\ln\Lambda^2$ 与常数 3 都属于积分常数（中含有限部分）。这些常数要与经典项 ${\cal F}=\frac12\tau_{cl}A^2$ 的归一化匹配才能唯一确定，原文省略了这一步；常数“3”也可直接对 $A^2\ln A^2$ 求二阶导（得 $2+1$）核对。]
\begin{equation}
{\cal F}_{1-loop}(A)=\frac{i}{2\pi} A^2 {\rm ln} \frac{A^2}{\Lambda^2}\,.
\label{oneloop}
\end{equation}
这里 $\Lambda$ 是一个固定的、动力学产生的标度，类似于 $\Lambda_{QCD}$。众所周知，由于 $N=2$ 超对称性，上述预势的单圈表达式不会收到更高阶的微扰修正。这与该理论中微扰 $\beta$ 函数只有单圈效应这一事实有关。于是，在该理论中，微扰 $\beta$ 函数的单圈性质与众所周知的反常的单圈性质是相容的。在 $N=1$ 理论中情况并非如此：那里的反常当然仍是单圈效应，但 $\beta$ 函数不是（例如见 \cite{SV1, SV2}）。

% [Tutor Note: “$N=2$ 的 $\beta$ 函数只有单圈”是 $N=2$ 非重整化定理的结论。系数可核对：在 $N=1$ 语言中单圈系数为 $3C_2(G)-T({\rm adj})$，对 $SU(N_c)$ 即 $3N_c-N_c=2N_c$，故 $\beta(g)=-g^3\cdot 2N_c/(16\pi^2)=-N_c g^3/8\pi^2$；$N_c=2$ 正给出原文的 $\beta(g)=-g^3/4\pi^2$。正因为如此，“反常严格单圈”与“预势在微扰论中精确到单圈”才能同时成立。对比 $N=1$：那里手征反常仍严格单圈，但 $\beta$ 函数会得到高阶修正（见 [36,37]），所以 $N=1$ 的预势没有这种单圈精确性。]

尽管 (\ref{oneloop}) 给出的预势在微扰论中是精确的，但如 \cite{Seiberg88} 所论证的，它确实会由于瞬子效应而收到非微扰修正。这些修正的一般形式可如下确定：首先，很清楚，来自瞬子数为 $k$ 的组态的 ${\cal F}$ 修正应正比于 $k$-瞬子因子 $exp(-8\pi^2k/g^2)$（因为预势是全纯函数，它不能从反瞬子组态获得修正）。利用该理论的单圈 $\beta$ 函数 $\beta(g)=-g^3/4\pi^2$，$k$-瞬子因子可以写成

% [Tutor Note: 单圈 $\beta$ 函数积分给出 $8\pi^2/g^2(\mu)=4\ln(\mu/\Lambda)$：由 $dg/d\ln\mu=-g^3/4\pi^2$ 得 $d(8\pi^2/g^2)/d\ln\mu=4$，于是 $e^{-8\pi^2k/g^2}=(\Lambda/\mu)^{4k}$；取 $\mu\sim a$（相关能标）即得下式，常数因子吸收进 $\Lambda$ 的定义。另：全纯性论证（见 [38,39]）排除反瞬子修正——反瞬子贡献依赖 $\bar g$ 而不全纯。]
$$
e^{-8\pi^2k/g^2} = \left(\frac{\Lambda}{a}\right)^{4k}\,.
$$
此外，仿照 Seiberg 在 \cite{Seiberg-n, Seiberg-poh} 中的做法，我们注意到：若赋予 $\Lambda$ 以荷 $2$，则被破缺的 $U(1)_{\cal R}$ 对称性便得到恢复。作了这一修改后，预势在 $U(1)_{\cal R}$ 下应作为一个荷为 $4$ 的场变换，且不带非齐次项。这意味着 $k$-瞬子修正也应正比于 $A^2$。把这些综合起来，包含一般非微扰修正的预势可以写成

% [Tutor Note: 给 $\Lambda$ 赋 $U(1)_{\cal R}$ 荷 2 的技巧值得说明：$\Lambda$ 是动力学标度而非基本场，但在反常意义下可以赋荷，使得 $\ln(A^2/\Lambda^2)$ 成为中性（$A$ 的 $U(1)_{\cal R}$ 荷为 2，故 $A^2$ 与 $\Lambda^2$ 的荷都是 4）；于是单圈项整体带荷 4，与 $k$-瞬子项 $(\Lambda/A)^{4k}A^2$ 的荷一致，$U(1)_{\cal R}$ 便以齐次方式恢复。这是 Seiberg 的全纯性论证（见 [38,39]）的关键一步。]
$$
{\cal F}=\frac{i}{2\pi} A^2 {\rm ln} \frac{A^2}{\Lambda^2}
+\sum_{k=1}^\infty\,{\cal F}_k\left(\frac{\Lambda}{A}\right)^{4k}A^2\,.
$$
系数 ${\cal F}_k$ 不依赖于场，这是因为在超对称理论中瞬子只通过零模对路径积分有贡献 \cite{DD}。系数 ${\cal F}_1$ 已在 \cite{Seiberg88} 中算出，结果非零。确定 ${\cal F}$ 的精确形式是 Seiberg 与 Witten 工作的主题。

${\cal F}$ 的单圈表达式也可以由积分 $\beta$ 函数得到，取 $A$ 作为积分限之一：在经典低能理论中——它只是把有质量场从微观拉格朗日量中丢掉（而不对它们作积分）——预势由 ${\cal F}(A)=\frac{1}{2} \tau_{cl}A^2$ 给出，其中 $\tau_{cl}=\frac{\theta}{2\pi}+i\frac{4\pi}{g^2}$。在量子理论中，对大 $a$（即 $A$ 的标量分量的真空期望值，vev），渐近自由起主导作用，理论处于弱耦合。因此在该极限下，把 $g$ 换成相应的跑动耦合 $g(a)$，就可以得到理论量子行为的一个良好近似（我们已经通过费米子场的手征重定义把 $\theta$ 取为零）。在大 $a$ 极限下，$g(a)$ 及 ${\cal F}(a)$ 可通过积分 $\beta$ 函数在微扰论中求得。

% [Tutor Note: “取 $A$ 为积分限”的细节：由 $\beta$ 函数积分 $1/g^2$ 得 $4\pi/g^2(a)=(2/\pi)\ln(a/\Lambda)+{\rm const}$，其虚部形式恰给出后文 $\operatorname{Im}\tau=\frac1\pi(\ln\frac{a^2}{\Lambda^2}+3)$ 的对数项；常数部分（含“3”）由与经典项 $\frac12\tau_{cl}A^2$ 的匹配确定。原文 “sompnent” 是 “component” 的拼写笔误（原文 PDF 第 50 页），“vev” 指 vacuum expectation value。]

对其他规范群，我们注意到对每个根 $\alpha$（$\alpha > 0$），都有有质量场 $W^{\alpha}, W^{- \alpha}$。对单链（simply-laced）群，$\alpha^{2} = 2$，且
$$
{\cal F} = \frac{i}{4\pi} \, \sum_{\alpha > 0} \, (\vec{\alpha} \cdot
\vec{a})^{2} \, {\rm ln}\, \frac{(\vec{\alpha}\cdot
\vec{a})^{2}}{\Lambda^{2}}\,.
$$
对数的系数来自单圈 $\beta$ 函数，它同时也导致未破缺的整体 $U(1)_{R}$ 对称性。

${\cal F}$ 的单圈表达式，加上“一个良定义的理论应当具有正定的动能项”这一要求，会引出一些很有意义的结论。对大的 $|a|$，利用 (\ref{oneloop}) 可以算出 $\tau(a)=\frac{i}{\pi}({\rm ln}\,\frac{a^2}{\Lambda^2}+ 3)$。这是一个多值函数，尽管由 ${\rm Im}\tau$ 给出的场空间度规是单值的。然而，由于 ${\rm Im}\tau(a)$ 是调和函数，它不可能有整体极小值。因此，若它被整体地定义，就不可能处处为正。所以，动能的正定性要求 $\tau(a)$ 只能局部地定义。这意味着，在复平面上 $\tau(a)$ 变为负值的区域，人们需要该理论的另一种描述。下一小节我们将说明如何得到这些等价描述。

% [Tutor Note: “调和函数不可能有整体极小值”用的是调和函数的最小值原理（maximum principle）：非常数调和函数在开区域内部不能取极值。这里有两层意思：${\rm Im}\,\tau$ 必须处处为正（否则度规不定、动能项变负），但单圈表达式 ${\rm Im}\,\tau=\frac1\pi(\ln\frac{a^2}{\Lambda^2}+3)$ 在小 $|a|$ 处会变负，故微扰描述只在弱耦合区（大 $|a|$）有效；$\tau(a)$ 的多值性（$\ln a^2$ 的对数支）正是后文单值变换（monodromy）的几何根源（原文 PDF 第 50 页）。]

\subsection{对偶}

为找出联系同一理论不同描述的对偶变换，我们考虑作用量玻色部分中的规范场项。在 Minkowski 空间中，采用约定 $(F_{\mu\nu})^2=-(\widetilde F_{\mu \nu})^2$ 与 $\widetilde{\widetilde F}=-F$，这些项可以写成
$$
\frac{1}{32 \pi} \, {\rm Im}\, \int \, \tau (a) (F+i\widetilde F)^2 =
\frac{1}{16 \pi} \, {\rm Im} \, \int \, \tau (a) (F^{2}+i\widetilde
FF) \,.
$$
现在我们把 $F$ 视为独立场，并通过引入一个拉格朗日乘子矢量场 $V_{D}$ 来实现 Bianchi 恒等式 $dF = 0$。为固定拉格朗日乘子项，把 $U(1)\subset SU(2)$ 作如下归一化：所有 $SO(3)$ 场都具有整数电荷。于是，如第 1 节所讨论的，$SU(2)$ 基本表示中的所有物质场将具有半整数电荷。在这一约定下，磁单极满足 $\epsilon^{0\mu\nu\rho}\partial_{\mu}F_{\nu\rho}=8\pi\delta^{(3)}(x)$。现在可以通过把 $V_D$ 与一个磁单极耦合来构造拉格朗日乘子项：

% [Tutor Note: 这里的 $8\pi$ 与 §1 的归一化约定严格绑定：把 $U(1)\subset SU(2)$ 归一化到“$SO(3)$ 场整数电荷”时，基本表示物质场的电荷为半整数，磁单极方程右端为 $8\pi\delta^{(3)}(x)$。若采用 $eg=2\pi$ 或 $eg=4\pi$ 等不同约定，$8\pi$、$1/8\pi$、$1/16\pi$ 等因子会整体改变；务必与 §1（1.2--1.4 小节）中电荷量子化与规范场归一化的约定对照核对（原文 PDF 第 51 页）。]
$$
\frac{1}{8\pi}\,\int\,V_{D \mu}\epsilon^{\mu\nu\rho\sigma}
\partial_{\nu}F_{\rho\sigma}=\frac{1}{8\pi}\,\int\,\widetilde F_{D}F = 
\frac{1}{16\pi}\,{\rm Re}\,\int\,(\widetilde
F_{D}-iF_{D})(F+i\widetilde F)\,,
$$
其中 $F_{D\mu\nu}=\partial_\mu V_{D\nu}-\partial_\nu V_{D\mu}$。把它加到规范场作用量上并对 $F$ 积分，给出对偶理论

% [Tutor Note: 中间略去的代数可补出：利用 2-形式的恒等式 $\widetilde F_D F=F_D\widetilde F$、$\widetilde F_D\widetilde F=-F_D F$ 及 $\operatorname{Re}(iX)=-\operatorname{Im}X$，可把 $\frac1{16\pi}\operatorname{Re}\int(\widetilde F_D-iF_D)(F+i\widetilde F)$ 化为 $\frac1{8\pi}\int\widetilde F_D F$；若把 $X=F+i\widetilde F$ 当作整体（注意 $X^2=2F^2+2iF\widetilde F$），对 $F$ 配方即得下一式的对偶作用量。原文以“Adding this ... and integrating over $F$”一句带过。]
$$
\frac{1}{32\pi}\,{\rm Im}\int \left(-\frac{1}{\tau}\right)
(F_D +i\widetilde F_D)^2=\frac{1}{16\pi}\,{\rm Im}\int\left(-\frac{1}{\tau}
\right) (F_D^2 +i\widetilde F_D F_D)\,.
$$
对偶化也可以在 $N=1$ 超对称语言中进行。此时作用量中的相关项为
$$
\frac{1}{8 \pi} \, {\rm Im} \, \int \, d^{2}\theta
\tau (A) W^{2}\,.
$$
Bianchi 恒等式被替换为 ${\rm Im}{\cal D}W = 0$。这可以通过引入一个矢量超场 $V_{D}$ 来实现，相应的拉格朗日乘子项变为
$$
\frac{1}{4\pi}\,{\rm Im}\,\int d^4x d^4\theta\, V_{D}{\cal D}W 
=\frac{1}{4 \pi}\,{\rm Re}\,\int d^4x d^4\theta\, i{\cal D}V_D W
=-\frac{1}{4\pi}\,{\rm Im}\,\int d^4x d^2\theta\, W_{D} W\,.
$$
把它加到作用量中并对 $W$ 积分，给出对偶作用量
$$
\frac{1}{8\pi}\,{\rm Im}\,\int d^2\theta \left(-\frac{1}{\tau(A)}
W^2_D \right)\,. 
$$
于是，对偶变换的效果是：把耦合到电荷的规范场替换为耦合到磁荷的对偶规范场，同时把规范耦合变换为
\begin{equation}
\tau\rightarrow\tau_D = -\,\frac{1}{\tau}\,.
\label{tautauD}
\end{equation}
我们认出这正是第 1 节的电磁对偶。另外注意，作用量在替换 $\tau\rightarrow\tau+1$ 下不变。这一变换与 (\ref{tautauD}) 中的变换一起生成 $SL(2,{\rm Z})$ 群，因此该群就是我们理论完整的对偶群。该群对 $\tau$ 的作用为

% [Tutor Note: $\tau\to\tau+1$ 对应 $\theta\to\theta+2\pi$：$\theta$ 项为 $\theta(32\pi^2)^{-1}\int F\widetilde F=\theta\nu$，在 $e^{iS}$ 中只依赖 $\theta\bmod 2\pi$，故作用量（路径积分）不变。$T:\tau\to\tau+1$ 与 (3.8) 的 $S:\tau\to-1/\tau$ 一起生成 $SL(2,{\bf Z})$；不同文献对生成元的取法（$S$ 或 $-S$、$T$ 或 $T^{-1}$）只差约定，核对时以下式 (3.11) 的矩阵为准。]
\begin{equation}
\tau\rightarrow{\frac{a\tau + b}{c\tau + d}}\,,
\label{tauSLtwoZ}
\end{equation}
其中 $ad-bc=1$ 且 $a,b,c,d \in {\rm Z}$。

$N=2$ 超对称性把规范场的对偶描述与标量场的对偶描述联系起来。为看出这一点，让我们引入 $h(A)=\partial{\cal F}/\partial A$。用它表示时，$\tau(A)=\partial h(A)/\partial A$，而标量动能项成为 ${\rm Im}\int d^4\theta\,h(A)\bar A$。对相应于 $(\ref{tautauD})$ 的对偶理论，让我们引入相应的对偶变量 $A_D,{\cal F}_D,h_D(A_D)$ 与 $\tau_D$。于是，方程 (\ref{tautauD}) 意味着 $A_D=h=\partial{\cal F}/\partial A$ 且 $h_D = -A$。在这一变换下，标量动能项保持其形式，

% [Tutor Note: $A_D=h$ 与 $h_D=-A$ 的推导：由 $\tau_D=\partial A_D/\partial A=-1/\tau=-\partial A/\partial A_D$，积分即得 $A_D=h(A)$、$h_D(A_D)=-A$（相差可加常数，原文取零）。这正是 $(A_D,A)$ 在 $S:\tau\to-1/\tau$ 下按 (3.11) 的 $S=\begin{pmatrix}0&1\\-1&0\end{pmatrix}$ 变换的含义；代入 $\operatorname{Im}\int d^4\theta\,h\bar A$ 即见其不变（$d^4\theta$ 积分不产生额外因子）。]
$$
{\rm Im}\int d^4\theta\,h(A)\bar A=
{\rm Im}\int d^4\theta\,h_D(A_D)\bar A_D\,. 
$$
下面我们把记号 $A_D$ 与 $h(A)$ 互换使用。

现在考虑完整的 $SL(2,Z)$ 群对 $A$ 与 ${\cal F}$，或等价地对 $A$ 与 $A_D=\partial{\cal F}/\partial A$ 的作用。变换 (\ref{tauSLtwoZ}) 意味着
\begin{equation}
\left(\begin{array}{c}A_D\\ A\end{array}\right)\rightarrow
\left(\begin{array}{cc} a & b \\ c & d \end{array}\right)
\,\left(\begin{array}{c}A_D\\ A \end{array}\right) \,.
\label{ASLtwoZ}
\end{equation} 
注意，一般地我们还可以在以上方程右端加上一个常数列矩阵 $C$。然而，当 $C$ 非零时，无物质情形下理论的 BPS 质量公式在上述变换下不再不变。因此在这种情况下应取 $C=0$。不过，在存在物质场时，BPS 质量公式会被修改，非零的 $C$ 是被允许的。这一情形将在最后一节讨论。在标量场空间上，变换 $\tau\rightarrow -1/\tau$ 与 $\tau\rightarrow \tau+1$ 由以下矩阵实现：
\begin{equation}  
S=\left(\begin{array}{cc}0&1\\-1&0\end{array}\right),\quad
T=\left(\begin{array}{cc}1&1\\0&1\end{array}\right).
\label{ST}
\end{equation}
这些矩阵生成 $SL(2,{\rm Z})$ 群。

${\cal F}$ 的变换可以直接由 (\ref{ASLtwoZ}) 得到，或等价地由下式得到：

% [Tutor Note: 原文 “equivanlently” 为 “equivalently” 的拼写笔误（原文 PDF 第 52 页）。]
\begin{eqnarray}
A'_D &=& a A_D + b A\,, \nonumber\\
A'   &=& c A_D + d A\,. \nonumber
\end{eqnarray}
利用第二式，可以把第一式对 $A'$ 积分，结果为

% [Tutor Note: 此处 $\alpha,\beta,\gamma,\delta$ 是 (3.10) 中 $SL(2,{\rm Z})$ 矩阵的四个元素（即 $\alpha=a,\ \beta=b,\ \gamma=c,\ \delta=d$），与本节其他地方的 $\alpha=g^2/4\pi$ 及 $U(1)_{\cal R}$ 转角无关——同一字母三种含义并存是原文记号的一大陷阱。可自行核对：要求 $\partial{\cal F}'/\partial A'=A'_D$ 就得到 $\beta\gamma=\alpha\delta-1$，即 $\alpha\delta-\beta\gamma=1$。]
$$
{\cal F}'=\frac{1}{2}\beta\delta A^2 +\frac{1}{2}\alpha\gamma A^2_D
+\beta\gamma A A_D + {\cal F}\,.
$$

现在回到 (\ref{metric}) 给出的模空间度规。用变量 $a_D$ 表示，它成为
$$
ds^2={\rm Im}\,da_D d\bar{a}=-\frac{i}{2}(da_D d\bar{a}-da d\bar{a}_D)\,,
$$
它在 $SL(2,{\rm Z})$ 变换下不变。为使描述更精确，注意模空间 ${\cal M}$ 是一维复流形，设 $u$ 是该流形上的全纯坐标。最终我们将把 $u$ 认同为 $\langle{\rm Tr}\phi^2\rangle$。设 $a$ 与 $a_D$ 是空间 $X \cong {\rm C}^2$ 上的坐标，在该空间上可以选取辛形式 $\omega={\rm Im}\,da_D\wedge d\bar{a}$。函数 $(a_{D}(u),a(u))$ 给出一个从 ${\cal M}$ 到 $X$ 的映射 $f$。换言之，它们确定了 $X$（视为 ${\cal M}$ 上的 $SL(2,{\rm Z})$ 丛）的一个截面。用坐标 $u$ 表示时，模空间上的度规具有如下形式

% [Tutor Note: 这段话把度规的 $SL(2,{\rm Z})$ 不变性几何化：$X\cong{\rm C}^2$ 上的辛形式 $\omega={\rm Im}\,da_D\wedge d\bar a$ 在 $SL(2,{\rm Z})$ 下不变（$\operatorname{Im}$ 部分对应 $S$ 保持的反对称形式），而模空间度规是 $\omega$ 沿截面 $f$ 的拉回，故自动不变。“截面”指 $(a_D(u),a(u))$ 作为 $X$ 上丛的局部选取；$u$ 只要求全纯，它与 $\langle{\rm Tr}\phi^2\rangle$ 的最终认同、以及小 $|a|$ 区域的解，将在后续小节确定（原文 PDF 第 52 页）。]
$$
ds^2 ={\rm Im}\,\frac{da_D}{du}\,\frac{d \bar{a}}{d\bar{u}}\,
dud\bar{u}=-\frac{i}{2}\left(\frac{da_{D}}{du}\,\frac{d\bar{a}}{d\bar{u}}
-\frac{d\bar{a}_D}{d\bar{u}}\,\frac{da}{du}\right)dud\bar{u}\,.
$$
这是与辛形式 $\omega$ 相联系的凯勒度规的拉回，因此明显具有 $SL(2,{\rm Z})$ 不变性。取 $u=a$，就回到原来的公式。注意，对任意的 $(a_D(u),a(u))$，该度规并非正定的。不过，我们后面得到的解会确定这些函数，使度规恒为正。

% [Tutor Note: “取 $u=a$ 回到原公式”是因为此时 $\frac{da_D}{du}=\tau$，故度规化为 $\operatorname{Im}\tau\,|da|^2$ 即 (3.5)。另注意：$SL(2,{\rm Z})$ 不变性只保证形式，不保证正定性；正是“动能正定”这一物理要求挑出正确的解（见后续小节）。]

\noindent\underline{高维情形：}
{}对秩为 $r$ 的群，度规具有如下形式

% [Tutor Note: 标题 “Higher Dimensions” 指模空间维数为 $r>1$（更高秩的规范群），不是时空维数；段首的 “{}” 是作者用来避免 \verb|\underline| 后紧跟文字时产生间距的技巧，这里照录。]
$$
ds^2 ={\rm Im}\,\frac{\partial^2{\cal F}}{\partial a_i\partial a_j}\, 
da_i d\bar{a}_j\,.
$$
引入一个空间 $X \simeq {\bf C}^{2r}$，坐标为 $a^{i}, a_{Dj}=\partial {\cal F}/\partial a^{j}$。给 $X$ 赋以辛形式 $\omega=\frac{i}{2}\,\sum_{i}(da^i\wedge d\bar{a}_{Di}-da_{Di} \wedge d\bar{a}^i)$ 与全纯 2-形式 $\omega_h =\sum_{i} \, da^{i} \wedge da_{Di}$。选取 $u^s\, (s=1,...,r)$ 作为模空间 ${\cal M}$ 上的坐标。于是 $a^i(u),a_{Dj}(u)$ 给出一个从 ${\cal M}$ 到 $X$ 的映射 $f$，使得 $f^{*}(\omega_{h}) = 0$（这一条件保证 $a_{Dj}=\partial {\cal F}/\partial a_{j}$）。因此，度规具有如下形式

% [Tutor Note: $f^*(\omega_h)=0$ 等价于 $\partial a_{Di}/\partial a^j=\partial a_{Dj}/\partial a^i$（反对称部分为零），这正是存在局部全纯函数 ${\cal F}(a)$ 使 $a_{Di}=\partial{\cal F}/\partial a^i$ 的积分条件（Poincaré 引理）。原文只说“这一条件保证…”，实际用到的是闭形式局部恰当这一标准事实。]
$$
ds^{2} = {\rm Im} \, \frac{\partial a_{Di}}{\partial u^{r}} \,
\frac{\partial \bar{a}^{i}}{\partial \bar{u}^{s}} \,
du^{r} d\bar{u}^{s}\,.
$$
这是与 $f^*(\omega)$ 相联系的度规，因此它在 $Sp(2r,{\rm R})$ 下不变：若记 $v^T =(a_D, a )$，则度规在 $v\rightarrow Mv + c$ 下不变，其中 $M \in Sp(2r,{\bf R})$。在无物质情形，$c=0$，而且只有 $Sp(2r,{\bf Z})$ 部分留存下来。

若模空间 ${\cal M}$ 具有非平凡结构，则当沿一条闭合回路绕行时，矢量 $v^T=(a_D, a)$ 会被单值群（monodromy group）的一个元素变换。正如后文将表明的，单值群是 $SL(2,{\rm Z})$ 的一个子群，确定它对于求解该问题是至关重要的。

% [Tutor Note: 单值变换（monodromy）的定位：把 $(a_D,a)$ 沿模空间中绕奇点的闭合回路作解析延拓，回到起点时一般不再是原来的值，而相差一个 $SL(2,{\bf Z})$ 变换；这正是 $\tau(a)$ 多值性（对数支）的全局体现，后文将用单值矩阵 $M_\infty$ 等拼接各个弱耦合描述。原文 “on being taken around a close loop” 中 “close loop” 应为 “closed loop”（拼写笔误，原文 PDF 第 53 页）。]

\subsection{双子质量}

如第 1、2 节所述，对微观的 $SU(2)$ 理论，BPS 界由 $M \geq \sqrt{2} |Z|$ 给出，其中，

% [Tutor Note: 原文 “microspcopic” 为 “microscopic” 的拼写笔误（原文 PDF 第 53 页）。此处的 $Z$ 是经典（树图）中心荷；下文 (3.12) 给出低能有效理论中的精确形式 $Z=a n_e+a_D n_m$。]
$$
Z=a(n_e + \tau_{cl}\, n_m)\,, 
$$
且 $\tau_{cl}=\frac{\theta}{2\pi}+i\,\frac{4\pi}{g^2}$。所有达到这一界的态都属于 $N=2$ 代数的短多重态。在更一般的 $N=2$ 理论中，例如我们一直在研究的低能有效作用量，该公式会略有修改。假设理论含有超多重态形式的物质场。当 $a\ne 0$ 时，这些在 $N=1$ 形式中由手征场 $M, \widetilde M$ 描述的场将变为有质量的。若一个超多重态带有电荷 $n_e$，则它与手征场 $A$ 的耦合由 $N=2$ 超对称性唯一确定为

% [Tutor Note: 这里的质量项可从 $N=2$ 的超势 ${\cal W}=\sqrt{2}\widetilde Q\Phi Q$（见 (3.2) 之前关于超势的讨论）得到：把 $\Phi$ 的标量分量取为真空期望值（原文在此处记作 $A$，通常记作 $a$），即得 $\sqrt{2}n_e a M\widetilde M$，于是 $a\ne0$ 时超多重态获得质量——这是希格斯机制在 $N=2$ 中的实现。]
$$
\sqrt{2}\,n_e \, A M \widetilde M\,.
$$
由此我们容易看出，对这样的态，$Z=a\,n_e$。另一方面，若考虑一个带有 $n_m$ 个单位磁荷的磁单极，则一个与第 1.6 小节非常相似的论证给出 BPS 界 $\sqrt{2}|n_m\,a_D|$，或等价地 $Z=n_m\,a_D$。一般地，可以对双子算出这一公式，得到
\begin{equation}
Z = a n_e + a_D n_m\,.
\label{ZaaD}
\end{equation}
若规范群为秩 $r$，则我们有
$$
Z = a^i n_{e,i} + h_i(a) n^i_m =a\cdot n_e +a_D \cdot n_m\,.
$$
由于质量是物理可观测量，质量公式在模空间上的单值变换下应当不变。因此，若 $v=(a_D, a)^T$ 变换为 $M v$（其中 $M\in Sp(2r,{\rm Z})$），则矢量 $w= (n_m, n_e)$ 应变换为 $w M^{-1}$。注意，这再次要求 $M$ 是整系数矩阵。

% [Tutor Note: $w=(n_m,n_e)$ 是行矢量，$Z=w^Tv$ 是不变量；等价的列矢量写法是按 $(M^T)^{-1}$ 变换。$M$ 为整系数矩阵是要求电荷格映射到自身，否则 $n_m,n_e$ 不再是整数——这就是“对偶群必须落在 $Sp(2r,{\bf Z})$ 中”的物理含义。]

在单值变换作用下，耦合矩阵 $\tau_{ij}=\partial a_Di/\partial a_j$ 变换为 $(A\tau+B)(c\tau+D)^{-1}$，其中 $M=\left(\begin{array}{cc}A&B\\C&D\end{array}\right)$。这与亏格 $r$ 的黎曼面在其模空间上单值群作用下的周期矩阵变换非常相似。这使得把我们的 $r$ 维真空模空间 ${\cal M}$ 认同为亏格 $r$ 黎曼面的模空间显得相当自然。变量 $a_i,a_{Dj}$ 于是与这些曲面的周期有关，并且可以计算出来。此外，${\rm Im}\,\tau$ 总是正定的，这就解决了出现负动能项的问题。为对我们的 $SU(2)$ 理论检验这一假设，我们首先必须确定模空间 ${\cal M}$ 上的单值结构。这是下一小节的主题。

% [Tutor Note: 原文此处有一处笔误（原文 PDF 第 54 页）：$\tau_{ij}$ 的变换写作 $(A\tau+B)(c\tau+D)^{-1}$，其中小写 $c$ 与矩阵块 $C$（$M=\begin{pmatrix}A&B\\C&D\end{pmatrix}$）不一致，正确形式应为 $(A\tau+B)(C\tau+D)^{-1}$。同页另有两处小笔误：“geting” 应为 “getting”；$\partial a_Di/\partial a_j$ 的下标应为 $\partial a_{Di}/\partial a_j$。公式与文字已按原文照录，此处如实指出。]
